\documentclass[11pt,a4paper]{article}

\usepackage[utf8]{inputenc}
\usepackage[T1]{fontenc}
\usepackage{lmodern}
\usepackage{amsmath,amssymb,amsthm}
\usepackage{graphicx}
\usepackage{hyperref}
\usepackage{authblk}
\usepackage{amsmath, amssymb, amsfonts, amsthm}
\usepackage{geometry}
\geometry{margin=3.5cm}
\usepackage{booktabs}
\usepackage{enumitem}

\title{The Emergence of Physical Reality within the Framework of the Projective Dynamic Logo}

\author{C\'edric Laubscher}
\affil{Independent researcher in theoretical and mathematical physics, Switzerland\\}

\date{\today}

\begin{document}

\maketitle

\begin{abstract}
We introduce the Projective Dynamic Logo (PDL), a foundational theoretical framework that seeks to reconstruct physical reality from purely relational axioms formulated on signed graphs, thereby obviating any prior assumption of an underlying space–time manifold or material substrate. Beginning from four primitive postulates—minimal binary pulsation, triangular coherence, minimal completeness, and logical optimisation—we derive a minimal closed configuration, namely the complete signed graph on four vertices and six edges, denoted (4,6). This structure is interpreted as a logically stationary regime, proposed to correspond to an electron at rest. 

On the basis of this construction, we develop a hierarchical combinatorial architecture representing the proton and show how several dimensionless or scale-defining constants—specifically the fine-structure constant $\alpha$, Boltzmann’s constant $k_B$, and an effective gravitational coupling—can be reformulated as ratios of internal coherence within the PDL framework. Finally, we outline a relational characterisation of proper time, a description of gravitation as spatial and temporal variation in coherence cost within an emergent metric structure, and a speculative cosmological scenario in which constraint events propagate through a pre-physical logical field.
\end{abstract}

\newpage
\section{Introduction}

Modern fundamental physics is grounded in a collection of highly successful theoretical frameworks, most prominently quantum field theory and general relativity. However, when these theories are extrapolated to extreme regimes—such as Planck-scale phenomena, strong-gravity environments, or the earliest epochs of cosmological evolution—their foundational concepts become increasingly difficult to formulate in a logically coherent and self-consistent manner. Core notions such as space, time, energy, and even that of a well-defined physical state lose their clear meaning, and the central challenge undergoes a shift from primarily empirical adequacy to conceptual and methodological coherence. In these domains, incremental refinements of existing models are no longer sufficient; instead, it becomes necessary to re-examine and possibly revise the very preconditions under which such models can be meaningfully constructed and applied.

A characteristic feature of many foundational programmes is that they presuppose, at the outset, either a pre-given space–time manifold, a particular dynamical law, or a probabilistic framework, and subsequently endeavour to reconstruct the empirically observed spectrum of particles and coupling constants. In contrast, the present approach begins from an almost minimal logical basis and investigates which structures are \emph{logically compelled} to arise if any entity is to exhibit temporally persistent existence. We do not postulate an underlying continuum, nor do we assume fields or material substances as primitive ontological ingredients. Instead, reality is modelled as a network of minimal logical closures, represented by finite signed graphs whose edges encode binary relations of agreement or disagreement, and whose configurations evolve according to a coherence-optimisation principle.

Within this Projective Dynamic Logo (PDL) framework, four fundamental axioms are posited: minimal binary pulsation, triangular coherence, minimal completeness, and logical optimisation. These axioms restrict the permissible organisation of relational structures such that closed, reproducible dynamical cycles can emerge. Under these constraints, the first admissible elementary stationary configuration is found to be a complete signed graph on four vertices and six edges, denoted \((4,6)\), for which there exists a sign assignment that renders all triangular subgraphs coherent, together with a global pulsation that preserves this coherence. This minimally closed configuration can be construed as a logical stationary regime that is naturally associated with an electron in a rest state.

The primary aim of this paper is to develop a relational framework and to examine the extent to which it can account, in a coherent and quantitatively non-trivial manner, for several salient features of microphysics and gravitation. After presenting the axioms and the underlying logical structure in a methodologically oriented way (Section~2), we demonstrate in Section~3 how the \((4,6)\) structure can be associated with the electron at rest, how a hierarchical combinatorial architecture for the proton can be derived from the same principles, and how several dimensionless constants—most notably the fine-structure constant \(\alpha\), Boltzmann's constant \(k_B\), and an effective gravitational coupling—can be reinterpreted as ratios quantifying internal coherence. Section~4 analyses the conceptual implications, strengths, and limitations of this approach, including its connections to quantum field theory, gravitation, and cosmology. Section~5 summarises the principal contributions and delineates directions for future research.

\section{Methods: Relational Axioms and Logical Framework}
\subsection{Minimal binary pulsation and informational entities}

Within the present theoretical framework, existence is not construed as being grounded in any pre-given substance, spatial manifold, temporal parameter, or energetic substrate. Rather, it is anchored in the possibility of sustaining a \emph{persistent distinction} within an initially pre-metric setting, i.e., a domain that is devoid of any a priori metric structure. Accordingly, the primary inquiry is not ``what exists?'', but instead ``under which minimal conditions can an entity be said to exist in a robust and reproducible manner?''. In a genuinely pre-metric context, any reference to location, duration, or intrinsic quantitative attributes is systematically excluded, as these concepts already presuppose an underlying geometric or metrical framework.

Under such constraints, an entity that manifests only once and never recurs is, in logical terms, indistinguishable from nonexistence: in the absence of repetition, no trace is available that could underwrite identification across multiple putative occurrences. Similarly, a strictly immutable state, devoid of any internal differentiation or external contrast, cannot be discriminated from a homogeneous, featureless background. Persistent existence therefore presupposes the emergence of a \emph{difference capable of reproducing itself}, such that it can be stably detected, tracked, and re-identified over time.

We formalise this requirement by positing a \emph{minimal binary pulsation}. Existence is conceived as an alternation between two mutually exclusive states, which we denote in an abstract manner as “agreement” and “disagreement”. This alternation does not occur within time in the usual metric sense: no period, frequency, or amplitude is assumed \emph{a priori}. It solely encodes the fact that “something occurs rather than nothing”, and that this “something” can recur in a logically coherent manner. At this level of abstraction, there is no notion of energy or signal; there is only the bare postulate of a reproducible alternation.

An elementary unit of such alternation is referred to as an \emph{informational entity}. It is characterised as “informational” not by virtue of conveying semantic content, but because it instantiates the most primitive form of reproducible distinction within a context lacking predefined metrics. Information, in this sense, designates the mere possibility of a difference that persists over time. 

Formally, we posit the existence of a discrete dynamical system acting on a finite set of binary variables (signs), such that at least one configuration belongs to a 2-cycle. The minimal invariant of existence is therefore identified with a logical oscillation between two complementary states, independently of any embedding in a continuous space–time framework.

\subsection{Triangular coherence and closed structures}

A single informational entity, even when characterized by a minimal binary oscillation, is insufficient to sustain a stable and determinate distinction in isolation. If existence is reduced to the mere occurrence of an alternation, two such entities remain indistinguishable in the absence of any relational structure: no internal criterion is available that would permit one to identify which entity is which. For a distinction to acquire operative validity, it must be defined \emph{relatively}, that is, with respect to something else. In this sense, relation does not function as an additional stratum superimposed upon already constituted entities; rather, it constitutes the necessary condition under which informational entities can be said to exist in a determinate and formally specifiable manner.

However, a single relation is insufficient to stabilise a reference. Any modification of one of the relata, or of the relation itself, immediately undermines the possibility of subsequent re-identification. A more articulated relational organisation is therefore required, one capable of sustaining reproducible distinctions without recourse to an external background structure. To formalise this requirement, we represent a logical configuration by a finite graph \(G = (V,E)\), where the vertices \(V\) correspond to informational entities and the edges \(E\) encode binary “agreement/disagreement” relations between them. Each edge \((i,j) \in E\) is assigned a sign \(s_{ij} \in \{+1,-1\}\), on which the minimal pulsation operates as a discrete dynamical process.

In this framework, \emph{triangles}—cycles of length three—serve as the fundamental motifs of logical closure. A triplet of vertices \((i,j,k)\) constitutes a triangle if the three edges \((i,j)\), \((j,k)\), and \((i,k)\) are elements of \(E\). Such a triangle is said to be \emph{coherent} if the product of the associated edge signs satisfies
\[
  s_{ij}\,s_{jk}\,s_{ik} = +1,
\]
That is, provided that the alternations around the cycle do not introduce an uncompensated discontinuity in the agreement/disagreement pattern. Conversely, a triangle is said to be \emph{violated} if at least one of its three edges is absent, or if the product of the edge signs equals \(-1\). Coherent triangles thus instantiate a minimal form of closure: the informational reference required to sustain the distinction between entities is fully contained within the cycle itself.

To quantitatively assess deviations from this closure condition, we introduce the quantity \(\varepsilon\), denoting the fraction of violated triangles in a given configuration. For a graph with \(n = |V|\) vertices, there exist \( \binom{n}{3} \) distinct unordered vertex triplets. We define
\[
  \varepsilon \;=\; \frac{\text{number of violated triangles in } G}{\binom{n}{3}},
\]
We adopt the convention \(\varepsilon = 0\) for all cases with \(n < 3\). Configurations for which \(\varepsilon\) is close to zero exhibit a high degree of triangular coherence and, consequently, a pronounced internal stability. In what follows, we characterise elementary closed structures by means of graphs in which coherent triangles constitute a connected network that does not depend on any external vertex or edge to maintain the distinctions they encode.

\subsection{Minimal completeness and the (4,6) structure}

The requirement of triangular coherence, taken in isolation, is insufficient to ensure that a relational configuration can sustain an autonomous and persistent distinction. A graph may exhibit internally coherent triangular substructures while nonetheless relying on an implicit external background for the stabilisation of reference. To rule out such configurations, we impose a \emph{minimal completeness} condition: all references necessary to define, reproduce, and maintain the distinction must be internal to the structure itself. Equivalently, an elementary closed structure must be logically self-sufficient and must not rely, even implicitly, on any external vertices or edges.

Formally, we consider finite graphs \(G = (V,E)\) in which each vertex is required to play an equivalent structural role. A configuration is said to be \emph{elementary stationary} if it satisfies the following conditions:
\begin{itemize}
  \item \(G\) is connected and complete: every pair of distinct vertices is joined by an edge.
  \item There exists an assignment of signs \(s_{ij} \in \{+1,-1\}\) to the edges such that all triangles are coherent, i.e.\ \(\varepsilon = 0\).
  \item The coherent triangles form a single connected component: the subgraph induced by the union of all coherent triangles is connected and contains at least two triangles sharing a common edge.
  \item All vertices are structurally equivalent: no vertex occupies a distinguished or privileged position, which encodes the absence of an externally imposed hierarchy.
  \item The discrete dynamics on the edge signs admits a binary pulsation \(s \leftrightarrow -s\) that preserves the coherence of all triangles.
\end{itemize}
An elementary stationary structure is termed \emph{minimal} if no proper subgraph of \(G\) satisfies all of the above conditions.

A systematic elimination of inadequate configurations demonstrates that no graph with fewer than four vertices can instantiate such a structure. In the case of one or two vertices, the formation of a triangle is impossible; with three vertices, at most a single triangle can be constructed, which is insufficient to produce a network of overlapping cycles capable of sustaining reference in a fully closed manner. Any attempt to design a non-hierarchical, closed configuration with fewer than four vertices either fails to attain structural completeness or necessarily introduces a latent asymmetry.

The first permissible instance occurs for \(n = 4\). The associated complete graph possesses
\[
  r_4 = \frac{4 \times 3}{2} = 6
\]
edges and is uniquely characterised, up to relabelling, by the pair \((4,6)\). Within this graph, there exists a particular assignment of signs such that all triangular faces are coherent (\(\varepsilon = 0\)); the set of coherent triangles forms a connected covering of the entire graph; and a global inversion of all signs induces a binary pulsation that preserves the coherence pattern. The \((4,6)\) configuration therefore furnishes the first explicit realisation of the minimal completeness requirement and will serve as the fundamental building block in the subsequent construction.

\subsection{Logical optimisation functional}

Even within the class of configurations that satisfy minimal completeness, a wide range of relational organisations remains logically admissible. It is therefore necessary to introduce a principle that discriminates which structures are expected to be dynamically sustainable, rather than merely logically or kinematically conceivable. Within the present framework, this selection is not enforced by an external dynamical law defined on a fixed background space–time. Instead, it is implemented through an \emph{internal optimisation principle}: among all admissible configurations, only those that realise an optimal trade-off between coherence, connectivity, and minimality are expected to be dynamically stable and hence to persist.

To formalise this notion, we characterise any given configuration by means of the following three quantities:
\begin{itemize}
  \item the fraction \(\varepsilon\) of violated triangles, as defined above;
  \item the normalised relational density \(\rho = r / r_{\max}(n)\), where \(r\) denotes the number of active edges and \(r_{\max}(n) = n(n-1)/2\) is the number of edges in the complete graph on \(n\) vertices;
  \item a minimality indicator \(m\), which assumes the value \(1\) if the configuration realises a closed structure without requiring any strictly richer substructure, and \(0\) otherwise.
\end{itemize}

We introduce a logical optimisation functional \(F(\varepsilon,\rho,m)\) subject to the following qualitative constraints:
\begin{itemize}
  \item \(F\) is strictly decreasing in \(\varepsilon\): configurations exhibiting fewer coherence violations are assigned lower values of the functional and are therefore preferred.
  \item \(F\) is (weakly) increasing in \(\rho\) over the relevant interval: an increase in relational density, up to the point of closure, is interpreted as a more efficient utilisation of the available internal relational budget.
  \item \(F\) attains a higher value when \(m=1\): elementary closed structures are favoured over non-minimal closed structures that can be decomposed into simpler closed components.
\end{itemize}
Operationally, we model \(F\) as
\[
  F(\varepsilon,\rho,m) \;=\; w_{\mathrm{coh}}\,f_{\mathrm{coh}}(\varepsilon) \;+\; w_{\mathrm{conn}}\,f_{\mathrm{conn}}(\rho) \;+\; w_{\mathrm{min}}\,m,
\]
where \(w_{\mathrm{coh}}, w_{\mathrm{conn}}, w_{\mathrm{min}} > 0\) are fixed weighting coefficients, and \(f_{\mathrm{coh}}, f_{\mathrm{conn}}\) are monotone functions encoding the qualitative behaviours specified above.

Within this framework, the \((4,6)\) configuration arises as the first local maximum of \(F\) at the fundamental level: it achieves \(\varepsilon = 0\), attains the maximal relational density \(\rho = 1\), and satisfies the minimality condition \(m = 1\). It is therefore designated as the elementary stationary configuration from which more complex hierarchical architectures will be systematically constructed in the subsequent sections.

\section{Results}

\subsection{The (4,6) structure as a prototype for the electron at rest}

Once the \((4,6)\) configuration has been identified as the first minimal elementary stationary structure, the subsequent step is to investigate its physical interpretation. Within the present theoretical framework, we do not \emph{a priori} identify this structure with the electron; instead, we inquire whether a natural correspondence can be established between its intrinsic logical dynamics and the empirically observed properties of an electron in its rest frame.

By construction, the \((4,6)\) graph realises a closed network composed of four informational entities connected by six signed relations, such that all triangular subgraphs are mutually coherent and a global sign inversion
\[
  s \;\longleftrightarrow\; -s
\]
induces a binary pulsation that preserves the overall coherence pattern. This internal 2-cycle defines an intrinsic periodic process that does not depend on any external time parameter. We interpret this logical pulsation as the abstract analogue of the electron’s Compton pulsation. More precisely, a physical realisation of the \((4,6)\) structure in the rest frame is associated with a characteristic frequency, and the reduced Planck constant \(\hbar\) is interpreted as the minimal quantum of action corresponding to a single coherence cycle of this internal dynamics.

In this framework, the rest energy of the electron is not treated as a fundamental input parameter, but rather as a quantitative measure of the coherence cost required to maintain the elementary stationary regime associated with the \((4,6)\) closure. The relation
\[
  E_0 \;=\; \hbar \,\omega_{\mathrm{C}}
\]
is accordingly reinterpreted as expressing a balance between the internal pulsation frequency \(\omega_{\mathrm{C}}\) of the underlying logical structure and the minimal action required per coherence cycle. Within this perspective, the electron at rest is understood as the first non-trivial, self-sustaining regime of existence admitted by the relational axioms, instead of being posited as an irreducible material particle defined on a pre-established space–time background.

\subsection{Hierarchical architecture of the proton}

\subsubsection{Valence cores and relational sea}

Whereas the electron is represented as a single elementary stationary closure of type \((4,6)\), the proton necessitates a more intricate, composite structural description. Empirical results from quantum chromodynamics (QCD) demonstrate that the proton cannot be adequately modeled as consisting of three pointlike quarks alone: the dominant contribution to its mass originates from an intensive, strongly non-perturbative gluon field, together with a fluctuating sea of virtual quark–antiquark pairs. Within the relational framework, this internal complexity is encoded through a hierarchical organisation that integrates multiple stationary domains embedded in, and dynamically coupled to, an extensive fluctuating background.

More precisely, the proton is described as a two-tier structure consisting of three local, quasi-stationary domains—hereafter referred to as the \emph{valence cores}—which are embedded within a dense \emph{relational sea}. Each valence core is represented as a superposition of logical tetrads and is specified by the integer parameters
\[
  n_u = 24, \qquad n_d = 28,
\]
for up-type and down-type cores, respectively. Collectively, the three cores account for
\[
  r_{\mathrm{valence}} = 930
\]
effectively stationary relational degrees of freedom. These are supplemented by a large ensemble of dynamical relations forming a collective sea, comprising
\[
  R_{\mathrm{sea}} \simeq 10^4
\]
active links, yielding a total relational content
\[
  R_{\mathrm{tot}} = R_{\mathrm{valence}} + R_{\mathrm{sea}} \approx 11\,017.
\]
This decomposition implies that approximately \(9\%\) of the relations are stationary, whereas about \(91\%\) are dynamical. This partitioning is qualitatively consistent with the standard QCD-based understanding, according to which only a small fraction of the proton’s mass is attributable to valence degrees of freedom, while the predominant contribution arises from gluonic fields and sea-quark components.

\subsubsection{Active surface and coupling to an electron-type (4,6)}

The hierarchical closure associated with the proton does not couple to an electron-type \((4,6)\) structure via its entire internal volume, but rather through a finite \emph{active surface}. This active surface is defined as the minimal relational envelope through which the proton can establish coherent interactions with an external elementary closure, such as a stationary electron.

Within the combinatorial framework, the active surface is characterised by an effective number of surface relations \(R_{\mathrm{surf}}\), obtained from the underlying valence content in conjunction with a simplified geometric optimisation argument. A natural prescription is to interpret approximately one third of the total valence relations as representing an effective core contribution, and to multiply this by a dimensionless geometric factor. This yields an estimate for the active surface of the order
\[
  R_{\mathrm{surf}} \approx 5.0 \times 10^2.
\]
This finite set of surface links constrains and channels the couplings between the proton’s internal hierarchical structure and an external \((4,6)\) electron configuration. In this perspective, the proton may be regarded as an “inverted mini black hole”: a densely organised internal structure whose effective external interactions are governed predominantly by a limited, information-carrying surface, rather than by the entirety of its internal volume.

\subsection{Coherence-Oriented Framework for the Interpretation of Fundamental Physical Constants}

\subsubsection{Fine-structure constant as a surface-to-volume coherence fraction}

Within the relational framework outlined above, the fine-structure constant \(\alpha\) can be reinterpreted as a \emph{surface-to-volume coherence fraction}. The total coherence budget associated with a composite entity such as the proton is primarily allocated to sustaining the internal hierarchical closure encoded by the valence cores and the surrounding relational sea. Only a limited fraction of this coherence is effectively projected onto the active surface and is therefore available for coupling to an external \((4,6)\) electron structure.

Let \(\mu = m_p/m_e\) denote the proton–electron mass ratio and \(r_{\mathrm{valence}} = 930\) the number of valence relations. An effective active surface can then be defined as
\[
  R_{\mathrm{surf}} \;=\; \frac{r_{\mathrm{valence}}}{3}\,\varphi,
\]
where \(\varphi\) is a dimensionless geometric optimisation factor. For \(r_{\mathrm{valence}}/3 \approx 310\) and \(\varphi \approx 1.618\), one obtains \(R_{\mathrm{surf}} \approx 5.0 \times 10^2\). This leads to an effective expression for the inverse fine-structure constant of the form
\[
  \alpha^{-1} \;\approx\; \mu\,R_{\mathrm{surf}}^{\,2},
\]
which numerically reproduces the empirical value \(\alpha^{-1} \simeq 137\) with a relative deviation of order \(10^{-4}\). In this interpretation, \(\alpha\) quantifies the fraction of the proton’s internal coherence that is accessible at its surface for electromagnetic coupling, rather than appearing as a primitive and unexplained fundamental parameter.

\subsubsection{A Proposed Rigorous Formulation for Boltzmann's Constant \texorpdfstring{$k_B$}{kB}}

In statistical physics, Boltzmann’s constant \(k_B\) plays a role analogous to that of \(\hbar\) in quantum mechanics: it fixes the scale that connects microscopic configurations to macroscopic thermodynamic observables such as temperature and entropy. From a relational standpoint, \(k_B\) should be interpreted as quantifying the minimal granularity of coherence associated with collective excitations in an ensemble of hierarchically organized structures, rather than being treated merely as a phenomenological proportionality constant.

Within the present theoretical framework, the same structural ingredients that characterize the proton architecture—namely, the number of relations in the \((4,6)\) configuration, the ratio \(\mu\), the decomposition into valence and sea components, and a small coherence–leakage parameter \(\varepsilon\)—can be combined into a parameter-free candidate expression for the Boltzmann constant:
\[
  k_B^{\mathrm{(PDL)}} \;=\; h\,\nu_0\,\frac{n_{\mathrm{rel}}}{\mu^2}\,\frac{R_{\mathrm{sea}}}{R_{\mathrm{tot}}}\,\frac{\varepsilon}{4},
\]
where \(h\) denotes Planck’s constant, \(\nu_0\) is the minimal pulsation frequency of the electron at rest, \(n_{\mathrm{rel}} = 6\) is the number of relations in the \((4,6)\) block, \(R_{\mathrm{sea}}\) and \(R_{\mathrm{tot}}\) are the numbers of sea and total relations in the proton, respectively, and \(\varepsilon\) represents the coherence–leakage fraction. Substitution of empirical values leads to
\[
  k_B^{\mathrm{(PDL)}} \;\approx\; 1.38 \times 10^{-23}\,\text{J}\,\text{K}^{-1},
\]
with a relative discrepancy of order \(10^{-2}\) with respect to the CODATA reference value. Although this identification remains speculative, it provides quantitative support for the conjecture that the thermal energy scale can be traced back to the same relational substratum that underlies mass, electric charge, and gravitation, without the introduction of additional continuous free parameters.

\subsubsection{Effective gravitational coupling from coherence leakage}

In conclusion, within this relational framework, gravitation may be reinterpreted as the metric manifestation of a minimal \emph{coherence leakage} from composite nucleonic structures into their relational environment. The central hypothesis is that the hierarchical closure defining a proton cannot be perfectly self-contained: a non-vanishing fraction \(\varepsilon\) of its coherence necessarily fails to achieve strictly local closure and must instead be absorbed or redistributed within the surrounding network of coexisting closures.

In regimes where a large number of nucleons are present, the individual leakage contributions accumulate and necessitate a reorganisation of the compatibility relations among neighbouring structural domains. At a coarse-grained (effective) level, this reorganisation can be represented as a deformation of an emergent metric, which is phenomenologically perceived as a gravitational field. By expressing the effective gravitational coupling \(G_{\mathrm{eff}}\) in terms of the same underlying parameters that characterise the proton–electron system—specifically \(\mu\), the propagation speed \(c\), the proton mass \(m_p\), the minimal action scale, and the leakage fraction \(\varepsilon\)—one obtains a value that is compatible with the empirically measured Newtonian gravitational constant, within quantifiable and controlled uncertainties.

In this framework, the relative weakness of gravitation is attributed to the small magnitude of coherence leakage compared to the total internal energy budget of nucleonic (or more generally, composite) structures. Its universality follows from the assumption that all massive bodies instantiate the same underlying hierarchical organization of closures. Consequently, gravitation is not postulated as a fundamental field defined on a pre-specified manifold; rather, it is construed as an emergent, large-scale manifestation of the systematic failure of composite logical structures to attain perfect closure, whereby their residual coherence is effectively redistributed into, and encoded as, a metric background.

\section{Discussion}

The foregoing results indicate that a nontrivially large subset of both microscopic and macroscopic physics can be reformulated in terms of purely relational structures defined on signed graphs. Starting from four axioms and a minimal notion of binary pulsation, we derived an elementary stationary structure \((4,6)\) that admits a natural interpretation as a logical prototype for the electron at rest, as well as a hierarchical architecture for the proton that qualitatively parallels the familiar QCD decomposition into valence and sea components. Within this setting, several dimensionless constants and scale-setting quantities — specifically \(\alpha\), \(k_B\), and an effective gravitational coupling — acquire a unified representation as ratios or manifestations of internal coherence.

The present framework also exhibits several clear limitations. First, many of the numerical correspondences are, at this stage, controlled approximations rather than high-precision predictions. The expressions obtained for \(\alpha\) and \(k_B\), although non-trivial and free of adjustable parameters in their formal structure, still display relative discrepancies at the \(10^{-4}\)–\(10^{-2}\) level and depend on specific identifications between integers and relational substructures whose physical and mathematical justification remains incomplete. Second, the treatment of gravitation as a form of coherence leakage and of cosmology as a sequence of constraint events in a pre-physical logical field is, by construction, highly schematic. A substantially more detailed comparison with general relativity and with observational data is required to evaluate the viability and empirical adequacy of this perspective.

At a conceptual level, the proposal challenges the conventional hierarchy in which space–time, fields, and matter are regarded as primitive, while information or coherence are considered derivative. In the present framework, coherence itself is elevated to the status of a fundamental invariant: physical entities are modelled as specific realisations of minimal logical closures, and standard quantities such as proper time, mass, and coupling constants are reinterpreted in terms of the costs and constraints associated with maintaining these closures within a network of coexistence. This situates the approach in relation to other informational and relational programmes in quantum foundations and gravitation, but with a distinctive emphasis on discrete combinatorial structures rather than continuum field configurations.

The current status of the PDL framework is best characterised as that of a \emph{developing} foundational proposal, rather than a fully established theory. Some results follow as relatively direct consequences of the axioms and can be further sharpened and generalised mathematically; others are conjectural, exploratory, or only loosely constrained by extant empirical data. Accordingly, the most productive way to assess its scientific relevance is not to view it as a finished alternative theory, but as a structured hypothesis that yields falsifiable numerical relations and concrete research questions. Systematically testing and refining these relations, improving the underlying combinatorial analysis, and confronting the emergent metric description with high-precision astrophysical and cosmological observations constitute natural directions for future work.

\section{Conclusions}

In this work, we have investigated the Projective Dynamic Logo (PDL) as a relational framework for the emergence of physical reality. Starting from four axioms formulated on finite signed graphs, we identified the complete \((4,6)\) structure as the first minimal elementary stationary closure and interpreted it as a logical prototype of an electron at rest. Building on this foundation, we proposed a hierarchical combinatorial architecture for the proton, composed of valence cores, a dense relational sea, and a finite active surface that mediates couplings to an electron-type \((4,6)\) structure.

Within this architecture, several quantities that are conventionally treated as independent empirical parameters—namely the fine-structure constant \(\alpha\), Boltzmann’s constant \(k_B\), and an effective gravitational coupling—can be reformulated as coherence-related ratios or as signatures of a small leakage of coherence into an emergent metric background. The resulting parameter-free expressions reproduce the correct orders of magnitude and give rise to non-trivial numerical relations that can, in principle, be sharpened or falsified by further analysis.

The PDL framework should therefore be viewed as an intermediate-level proposal: sufficiently well-defined to yield explicit combinatorial models and quantitative constraints, yet still flexible and incomplete. Future research will focus on improving the mathematical characterisation of admissible closures, refining the proton architecture, extending the framework to other particles and bound systems, and developing a more detailed comparison with quantum field theory, general relativity, and cosmological observations.

\newpage

\begin{thebibliography}{99}

\bibitem{Laubscher_PDL_Intro}
C.~Laubscher,
\newblock ``Introduction to the Projective Dynamic Logo (PDL),''
\newblock Zenodo (2026), \href{https://doi.org/10.5281/zenodo.18463130}{doi:10.5281/zenodo.18463130}.

\bibitem{Laubscher_PDL_Long}
C.~Laubscher,
\newblock ``The Emergence of Physical Reality within the Framework of the Projective Dynamic Logo (PDL),''
\newblock Zenodo (2026), \href{https://doi.org/10.5281/zenodo.18462686}{doi:10.5281/zenodo.18462686}.

\bibitem{Laubscher_LDP_FR}
C.~Laubscher,
\newblock ``L'\'emergence de la r\'ealit\'e physique dans le cadre du Logo Dynamique Projectif (LDP),''
\newblock Zenodo (2026), \href{https://doi.org/10.5281/zenodo.18475542}{doi:10.5281/zenodo.18475542}.


\bibitem{QCD_Review}
C.~Amsler et al.,
\newblock ``Proton structure and quantum chromodynamics,''
\newblock {\em Physics Reports}, vol.~\dots, pp.~\dots, year~\dots.

\bibitem{GR_Textbook}
S.~Weinberg,
\newblock {\em Gravitation and Cosmology},
\newblock Wiley, 1972.

\bibitem{Relational_Approaches}
C.~Rovelli,
\newblock ``Relational quantum mechanics,''
\newblock {\em International Journal of Theoretical Physics}, 35, 1637--1678, 1996.

\bibitem{BH_Entropy}
J.~D. Bekenstein,
\newblock ``Black holes and entropy,''
\newblock {\em Physical Review D}, 7, 2333--2346, 1973.

\bibitem{Hawking_Radiation}
S.~W. Hawking,
\newblock ``Particle creation by black holes,''
\newblock {\em Communications in Mathematical Physics}, 43, 199--220, 1975.

% Ajoute ici toutes les autres références pertinentes :
% - articles sur la constante de structure fine
% - textes sur les approches informationnelles / relationnelles
% - références QCD standard, cosmologie, etc.

\end{thebibliography}


\end{document}
